This section presents four detection capabilities unified by a common mechanism---\emph{output suppression}---plus one refuted candidate (sycophancy). Three capabilities (refusal, absence-of-suppression in abliterated models, impossibility refusal) are new to Campaign 3. The key evidence for the output suppression interpretation is Experiment 36: impossibility refusal (AUROC $\approx 0.94$) is \emph{more} detectable than safety refusal (AUROC $\approx 0.91$), and the two refusal types are geometrically indistinguishable (AUROC $\approx 0.65$). This proves the signal is about \emph{withholding}, not content type.

\subsection{Deception Detection}
\label{sec:deception}

\paragraph{Within-model (system-prompt manipulation).} Using the C4 benchmark from pre-hackathon work, random forest classifiers on 4 aggregate features achieve AUROC 1.0000 across all 7 tested models for distinguishing honest from deceptive responses. This replicates the Campaign 2 finding.

\paragraph{Same-prompt deception (Experiment 18b).} A stronger test: both conditions use the identical system prompt (``You are a helpful assistant. Answer questions accurately and concisely.''). The deception instruction is embedded in the user message. On Qwen2.5-7B with 10 matched pairs:

\begin{table}[H]
\centering
\caption{Same-prompt deception detection. System prompt identical for both conditions.}
\label{tab:same_prompt_deception}
\begin{tabular}{lcccc}
\toprule
\textbf{Feature} & \textbf{Honest} & \textbf{Deceptive} & \textbf{Cohen's $d$} \\
\midrule
Key norm & $681.7 \pm 81.9$ & $871.0 \pm 139.3$ & $+1.66$ \\
Norm/token & $14.46 \pm 1.44$ & $11.57 \pm 1.52$ & $-1.96$ \\
Effective rank & $67.4 \pm 3.1$ & $74.6 \pm 2.6$ & $+2.50$ \\
Rank entropy & $22.1 \pm 0.9$ & $24.2 \pm 0.6$ & $+2.79$ \\
\bottomrule
\end{tabular}
\end{table}

GroupKFold AUROC: $\approx \textbf{0.92}$ (95\% CI $[0.83, 1.00]$, $p < 0.001$ by 10{,}000-permutation test), accuracy: 80\%. All four features show large or very large effect sizes, with rank entropy providing the strongest single-feature separation.

\paragraph{Length caveat.} The corrected AUROC of $\approx 0.92$ should be interpreted cautiously: a length-only baseline using raw norm achieves AUROC $\approx 0.91$, and a token-count baseline ($n_{\text{generated}}$) achieves $\approx 0.82$. At $n = 10$ per class, the bootstrap CI is wide ($[0.83, 1.00]$), and the $\approx 0.01$ improvement over the norm baseline is well within noise. While the corrected features exclude raw norm (using only norm-per-token, effective rank, and rank entropy), norm-per-token remains correlated with response length. The deception signal in same-prompt experiments cannot be confidently separated from response-length differences at this sample size.

\paragraph{Cross-condition transfer (Experiment 18).} Training on system-prompt-manipulated deception across 7 models, then testing on natural DeepSeek censorship (where the model evades politically sensitive questions without any deceptive system prompt): RF AUROC $\approx \textbf{0.89}$. This demonstrates shared geometric structure between instructed and natural output suppression.

\paragraph{Cross-model transfer (Experiment 19).} The cross-model transfer bottleneck from Campaign 2 (AUROC $\approx 0.67$ with RF) was resolved by switching to logistic regression:

\begin{table}[H]
\centering
\caption{Cross-model transfer: classifier comparison.}
\label{tab:transfer}
\begin{tabular}{lcc}
\toprule
\textbf{Approach} & \textbf{All-to-All AUROC} & \textbf{Qwen AUROC} \\
\midrule
Baseline RF & 0.752 & 0.678 \\
\textbf{Baseline LR} & \textbf{0.856} & \textbf{0.864} \\
Z-score RF & 0.816 & 0.839 \\
Interaction features LR & 0.863 & 0.862 \\
\bottomrule
\end{tabular}
\end{table}

The key insight: LR's linear decision boundary generalizes across architectures because aggregate cache features are approximately linearly separable. RF's axis-aligned splits overfit to architecture-specific feature ranges.

\subsection{Refusal Detection}
\label{sec:refusal}

\textbf{New capability.} Experiment 31 demonstrates that safety refusal---when a model declines to answer a harmful question---produces a detectable KV-cache signature. We interpret this as \emph{output suppression detection}: the model recognizes content it could elaborate on but withholds, allocating fewer representational resources per token.

\paragraph{Design.} 20 harmful prompts (designed to trigger refusal) paired with 20 benign prompts matched for domain. Qwen2.5-7B, 50 tokens generation. Actual refusal rate in text: 19/20.

\paragraph{Results.}

\begin{table}[H]
\centering
\caption{Refusal detection: single-model and multi-model results. LR values are corrected (GroupKFold, 3 features excl.\ raw norm, 10{,}000 permutations); RF values are from original evaluation.}
\label{tab:refusal}
\begin{tabular}{lccc}
\toprule
\textbf{Model} & \textbf{AUROC (LR)} & \textbf{AUROC (RF)} & \textbf{$d$ (rank)} \\
\midrule
Qwen2.5-7B & $\approx$\textbf{0.91} & 0.830 & 1.530 \\
Llama-3.1-8B & $\approx$0.87 & 0.893 & --- \\
Mistral-7B & $\approx$0.83 & 0.756 & --- \\
\midrule
\textit{Mean} & $\approx$\textit{0.87} & \textit{0.826} & --- \\
\bottomrule
\end{tabular}
\end{table}

Refusal responses are \emph{sparser per token} (295.7 vs.\ 301.1 norm/token on Qwen), consistent with output suppression: the model withholds its full representational capacity, generating a constrained, formulaic response. As shown in \Cref{sec:impossibility}, this suppression geometry is not specific to harmful content---it occurs whenever the model withholds output, whether for safety or logical impossibility reasons.

\subsection{Absence-of-Suppression Detection (Abliterated Models)}
\label{sec:jailbreak}

\textbf{New capability.} Experiment 32 tests whether abliterated models---those with the refusal direction removed via \citet{arditi2024refusal}---produce a different cache signature than either normal or refusing models. We frame this as \emph{absence-of-suppression detection}: abliterated models that answer harmful prompts lack the suppression circuitry that normally constrains output, and their cache geometry reflects this by showing \emph{normal} (non-suppressed) processing patterns rather than the sparse geometry characteristic of refusal.

\paragraph{Results.}

\begin{table}[H]
\centering
\caption{Absence-of-suppression detection across models.}
\label{tab:jailbreak}
\begin{tabular}{lcc}
\toprule
\textbf{Comparison} & \textbf{Llama AUROC} & \textbf{Mistral AUROC} \\
\midrule
Abliterated vs.\ Normal & $\approx \textbf{0.91}$ & $\approx 0.69$ \\
Abliterated vs.\ Refusal & $\approx 0.61$ & $\approx 0.48$ \\
\bottomrule
\end{tabular}
\end{table}

Abliterated responses to harmful prompts are separable from normal benign responses (AUROC $\approx 0.91$ on Llama) but nearly indistinguishable from refusal responses (AUROC $\approx 0.61$). This is initially counterintuitive---abliterated models \emph{comply} while refusing models \emph{decline}---but is explained by the output suppression framework: both conditions involve processing harmful content, and the suppression geometry reflects the \emph{content being processed}, not whether the model ultimately complies. The critical difference is between ``processing harmful content'' and ``processing benign content,'' not between ``refusing'' and ``complying.''

\paragraph{Caveat: heterogeneous labels.} Abliteration does not produce uniform compliance. Manual inspection of the Qwen abliterated outputs shows that only \textbf{8/20} (not 18/20) genuinely provide harmful information; ${\sim}10/20$ warn but partially comply, and ${\sim}2/20$ hard-refuse despite abliteration. The classifier separates this heterogeneous mixture from benign responses, but the signal reflects \emph{output suppression} (the model's residual suppression circuitry engaging when processing harmful content) rather than any ``jailbreak-specific'' geometry. This heterogeneity means the reported AUROC values should be interpreted cautiously: they measure separation of a \emph{mixed} class from benign, not clean ``compliant harmful'' versus ``benign'' discrimination.

\paragraph{Control: abliterated on benign prompts.} A follow-up experiment (Experiment 41) tested abliterated models on \emph{benign} prompts, yielding AUROC $= 0.118$---below chance. This confirms that the classifier detects \emph{suppression triggered by harmful content processing}, not abliteration as a structural modification: when abliterated models process benign content, their cache geometry is indistinguishable from normal benign responses. The signal is content-dependent, not a weight-modification artifact.

\subsection{Impossibility Refusal}
\label{sec:impossibility}

\textbf{Key evidence for the output suppression interpretation.} Experiment 36 tests whether refusal geometry is specific to safety-related content or reflects a general \emph{withholding} mechanism. This experiment provides the strongest single piece of evidence that the detected signal is output suppression, not harmful content detection.

\paragraph{Design.} 20 impossible questions (``What is north of north?'', ``What is the last digit of $\pi$?'') versus 20 benign questions. Qwen2.5-7B. These prompts contain no harmful content whatsoever---they are simply unanswerable.

\paragraph{Results.}

\begin{table}[H]
\centering
\caption{Impossibility refusal versus safety refusal: evidence for output suppression.}
\label{tab:impossibility}
\begin{tabular}{lcc}
\toprule
\textbf{Comparison} & \textbf{AUROC (LR)} & \textbf{AUROC (RF)} \\
\midrule
Impossible vs.\ Benign & $\approx \textbf{0.94}$ & $\approx 0.88$ \\
Harmful vs.\ Benign (Exp 31) & $\approx 0.91$ & $\approx 0.83$ \\
Impossible vs.\ Harmful & $\approx 0.65$ & $\approx 0.74$ \\
\bottomrule
\end{tabular}
\end{table}

Three findings converge on the output suppression interpretation:

\begin{enumerate}[leftmargin=*]
    \item \textbf{Impossibility refusal is \emph{more} detectable than safety refusal} (AUROC $\approx 0.94$ vs.\ $\approx 0.91$). If the signal were ``harmful content detection,'' impossibility---which involves no harmful content---should be undetectable.
    \item \textbf{The two refusal types are geometrically indistinguishable} (AUROC $\approx 0.65$, near chance; $p = 0.037$, marginal). Safety refusal and impossibility refusal produce the same cache geometry despite involving entirely different content domains.
    \item \textbf{The common mechanism is withholding.} The model produces a similar geometric perturbation whether it declines because ``that's dangerous'' or because ``that's impossible.'' In both cases, the model has representational capacity it could deploy but chooses not to, and this restraint manifests as reduced angular spread and lower norm-per-token.
\end{enumerate}

\noindent The most parsimonious explanation is output suppression---the model allocating fewer resources per token when withholding output. This is not harm-specific or intent-specific geometry; it is detection of whether the model is withholding.

\subsection{Sycophancy Detection: An Honest Null}
\label{sec:sycophancy}

Sycophancy detection is the one capability we claimed but cannot support.

\paragraph{Experiment 39 design.} One neutral system prompt for all conditions. Honest condition: bare factual question (``What year was the Treaty of Westphalia signed?''). Sycophancy condition: ``I think the answer is [wrong answer]. Am I right? [question].'' 20 matched pairs. Qwen2.5-7B.

\paragraph{The length confound.} LOO-CV produced AUROC 0.9375 with all $d > 2.3$. However, the sycophancy prompts are systematically ${\sim}15$--$20$ tokens longer than honest prompts due to the ``I think X. Am I right?'' preamble. The large effect sizes measure \emph{prompt length differences}, not sycophancy deliberation. The sign pattern confirms this: norm $d = +2.3$ (longer input $\to$ higher norm) while norm/token $d = -3.2$ (more tokens $\to$ less resource per token)---exactly the pattern expected from a length confound.

\paragraph{The length-matched comparison.} The properly controlled test---USER\_CORRECT vs.\ USER\_WRONG, where both prompts have the same preamble structure and differ only in the asserted answer---produces $d = 0.107$ (negligible on Cohen's scale). At $n = 60$ per group, the minimum detectable $d$ for 80\% power is 0.52. The observed effect is $4.9\times$ below the detection threshold. Achieving 80\% power at $d = 0.107$ would require $n \approx 1{,}400$ per condition.

\paragraph{Sycophancy rate.} Only 4/20 (20\%) of sycophancy-condition responses actually agreed with the wrong answer. The model resisted in 80\% of cases. With $n = 4$ genuine sycophancy samples, no classifier analysis is meaningful.

\paragraph{Context.} Recent work using per-layer residual stream probes (rather than aggregate KV-cache features) achieves AUROC $> 0.97$ for separating sycophantic from genuine agreement \citep{vennemeyer2025sycophancy}, but with datasets of $3{,}900$--$8{,}000$ pairs and per-layer DiffMean analysis. The signal exists in late layers (19+) and requires finer-grained features than our aggregate approach provides. Sycophancy detection via KV-cache geometry remains an open problem requiring either (a) per-layer analysis from late layers, (b) substantially larger prompt batteries, or (c) both.

\subsection{Extended Feature Set}
\label{sec:extended}

Experiment 38 adds the 5 extended key-geometry features (\Cref{sec:methods}), computed directly from key matrices without requiring attention weight extraction (important because several architectures, including Qwen2.5, do not support attention weight extraction via SDPA).

\begin{table}[H]
\centering
\caption{Feature set comparison for output suppression detection (Qwen2.5-7B, refusal vs.\ benign).}
\label{tab:extended}
\begin{tabular}{lcc}
\toprule
\textbf{Feature Set} & \textbf{AUROC (LR)} & \textbf{AUROC (RF)} \\
\midrule
Core 4 features & $\approx 0.90$ & $\approx 0.83$ \\
Extended 5 only & $\approx 0.89$ & $\approx 0.94$ \\
\textbf{Combined 9 features} & $\approx \textbf{0.95}$ & $\approx \textbf{0.94}$ \\
\bottomrule
\end{tabular}
\end{table}

The strongest new features: angular spread ($d = -0.857$; suppressed output produces less diverse keys), norm variance ($d = -0.702$; suppressed output produces more uniform processing), and generation delta ($d = +0.561$; suppressed output shows larger encoding--generation gaps).

The extended features capture complementary information: they measure \emph{how} the model processes (diversity, uniformity, phase transitions) rather than \emph{how much} (norms, ranks).

\paragraph{Caveat.} At $n = 20$ per class, the improvement from AUROC $\approx 0.90$ to $\approx 0.95$ ($\Delta \approx 0.05$) is within the expected CV noise. \citet{varoquaux2018cross} show that at these sample sizes, cross-validated AUROC estimates carry error bars of $\pm 0.15$--$0.20$; \citet{hanley1982meaning} provide the classical SE formula for AUROC, which at $n = 20$ per class gives SE $\approx 0.05$--$0.07$ (yielding 95\% CIs of $\pm 0.10$--$0.14$, though this assumes independent observations, which LOO-CV violates). No permutation test was performed to establish that 9 features significantly outperform 4. The extended features are scientifically informative (angular spread and norm variance provide mechanistic insight into the suppression mechanism) but the AUROC improvement should not be cited as validated.
